%!TeX root=csp.tex
\section{Defects in the sources}\label{sec:defects}

A blueprint earns its keep by finding the places where the source cannot be
transcribed. This section collects them. They are not criticisms of
\cite{ZhukSimplified}, which is unusually careful prose --- across nine thousand
lines of \TeX{} it says ``obvious'' three times and ``clearly'' never. They are
the residue that any informal exposition leaves, and each one is a decision a
formalization has to make before it can write the corresponding statement.

Items marked \textbf{blocking} must be resolved before the affected statement
can be written in Lean at all. Items marked \textbf{convention} are places
where the source is consistent but leaves the reader to supply a reading.

\subsection{Blocking}

\begin{enumerate}\tightlist

\item[\textbf{D1}] \textbf{Lemma~\ref{lem:ba-center-implies} is false as
  printed.} \cite[Lemma~19]{ZhukSimplified} drops two hypotheses present in its
  cited source: subdirectness of $R$ in the first coordinate, and a common
  absorbing term in the $\TBA$ case. Counterexample and repair in
  Hazard~\ref{haz:ba-center-implies}. Used three times in \S3 and three times
  in \S4 of the source.

\item[\textbf{D2}] \textbf{Corollary~\ref{cor:prop-quotient} is never proved.}
  It is stated in \S2.3, used four times inside \S5 and once in \S3, and
  neither proved nor restated anywhere. It is Lemma~\ref{lem:propagation}
  specialised to the canonical surjection $\mathbf A\to\mathbf A/\delta$, but
  the source does not say so, and the specialisation is not quite mechanical:
  clause (m)'s hypothesis ``$B/\delta$ is $\TS$-free'' has to be matched against
  clause (fm)'s ``$f(B)$ is $\TS$-free''.

\item[\textbf{D3}] \textbf{There is a citation cycle in \S5.} Tracing the
  proofs: Lemma~9 $\Rightarrow$ Lemma~7 $\Rightarrow$ Lemma~86 $\Rightarrow$
  Lemma~85 $\Rightarrow$ Corollary~22 $\Rightarrow$ Theorem~21 $\Rightarrow$
  Lemma~8 $\Rightarrow$ Lemma~9, closed where Lemma~85's $T=\TC$ case invokes
  Corollary~22. Nothing in \S5.4 or \S5.7 can be formalized until the cycle is
  broken. The break appears to require a mixed essentiality lemma for three
  distinct central subuniverses, which is not in the paper.

\item[\textbf{D4}] \textbf{Theorem~\ref{thm:stable-intersection}(c) asserts
  $m=2$ without argument.} The proof reduces to pairs and shows the types agree;
  for $\TPC$, $\sigma_{1}=\sigma_{2}$ forces $m=2$. Nothing rules out $m\ge 3$
  with all types $\TC$. A Lean statement copying the source is unprovable as it
  stands. Same fix as D3.

\item[\textbf{D5}] \textbf{Theorem~\ref{thm:codim-one} hypothesis (ii) is too
  weak for its proof.} The source says ``$\mathbf D_{x_{i}}$ is $\TS$-free'';
  the proof, and the informal claim it formalises, need ``no nontrivial binary
  absorbing or central subuniverse on $\mathbf D_{x_{i}}$'' --- which is the
  \emph{commented-out} gloss sitting immediately below the hypothesis in the
  source. $\TS$-freeness is strictly weaker (Hazard~\ref{haz:S-type}).
  Strengthen the hypothesis.

\item[\textbf{D6}] \textbf{Lemma~\ref{lem:connected}(a) applies
  Lemma~\ref{lem:bridge-from-relation} without one of its hypotheses.} The
  missing hypothesis can fail; a witness is
  $R=\{(x,z,y)\in\mathbb Z_{4}\times\mathbb Z_{2}\times\mathbb Z_{4}:
  x\equiv z\equiv y \pmod 2\}$. In \cite{ZhukJACM} the needed tuples came from
  \emph{criticality}, a notion \cite{ZhukSimplified} dropped. The fix is to add
  ``every constraint relation is critical'' to the definition of a connected
  instance --- available, since crucial implies critical --- at the cost of
  reimporting about a page of the older machinery.

\item[\textbf{D7}] \textbf{Two gaps in Theorem~\ref{thm:main-inductive}.}
  In Case~1 the inductive hypothesis yields (1c) relative to $D^{(2)}$ while the
  goal is (1c) relative to $D^{(1)}$; the source writes only ``we derive the
  required conditions''. And both endgames conclude ``(1b) if the solution set
  is subdirect, (1c) otherwise'', but (1c) demands a \emph{linked} connected
  subinstance where only connectedness was established --- linkedness has to be
  extracted from irreducibility, which needs a fragmentation descent and a
  separate treatment of the empty-solution-set case.

\item[\textbf{D8}] \textbf{A lemma the proof needs is commented out of the
  source.} \texttt{LEMMinimalContainingIsMinimal} --- that an inclusion-minimal
  $\MT T$-subset containing a given element is minimal among $\MT T$-subsets ---
  is commented out along with its citation, yet is genuinely used in Case~2 of
  Theorem~\ref{thm:main-inductive}. It is true and must be reproved.

\item[\textbf{D9}] \textbf{The algorithm's pseudocode does not match the theorem
  it invokes.} \textsc{SolveLinear} \emph{removes} constraints, whereas
  Theorem~\ref{thm:codim-one}(vii) is about \emph{weakening} them; and
  primitive \textup{(p2)} returns a union of affine subspaces, which the next
  line treats as affine. \cite{ZhukSimplified} labels its pseudocode a sketch
  and points to \cite{ZhukJACM} for the precise algorithm, so this is a defect
  of the sketch rather than of the proof --- but target \textup{(T1)} is a
  statement about the algorithm, so it must be repaired from \cite{ZhukJACM}
  before \textup{(T1)} can be stated.

\item[\textbf{D10}] \textbf{The arity in Lemma~\ref{lem:special-wnu} is suspect.} The
  imported statement gives a special WNU of arity $n^{n!}$; the period of
  $y\mapsto w(x,\dots,x,y)$ need not divide $n^{n!}$ --- take $n=3$ and an
  element of order $5$. The remedy is to weaken the statement to ``\emph{some}
  arity $N$'', which is all the dichotomy proof uses. The precise arity matters
  only for \S4 of the source, which is cut.

\end{enumerate}

\subsection{Convention}

\begin{enumerate}\tightlist

\item[\textbf{C1}] $\mathbf Z_{p}\in\Vn$ requires $p\mid n-1$
  (Hazard~\ref{haz:Zp}).
\item[\textbf{C2}] $\cover\sigma$ is a tolerance, not a congruence
  (Hazard~\ref{haz:irreducible}). Typing it as a congruence collapses the
  linear/PC distinction silently.
\item[\textbf{C3}] The empty set is a subuniverse of every idempotent algebra,
  and is vacuously absorbing and vacuously central. Read syntactically, the
  clause defining $\subt{\TS}$ is therefore universally true. The global
  convention ``we do not allow empty subuniverses'' repairs it, but the clause
  does not inherit that convention on its own
  (Convention~\ref{conv:empty}).
\item[\textbf{C4}] Whether $\sigma=A^{2}$ counts as irreducible is genuinely
  ambiguous --- it turns on whether the empty intersection is allowed --- and
  $\cover\sigma$ does not exist in that case. \cite{ZhukSimplified} drops the
  word ``proper'' that \cite{ZhukJACM} has. Legislating that the empty family
  \emph{is} allowed makes irreducibility imply properness, which is the
  convenient choice and is what the Lean development adopts.
\item[\textbf{C5}] $\lll$ is data, not merely a proposition: several \S5 proofs
  speak of ``the congruences coming from $C\lll A$'' and of minimal chain
  length (Hazard~\ref{haz:lll-chain}).
\item[\textbf{C6}] ``Dimension'' of $\prod\mathbb Z_{q_{i}}$ with distinct
  primes is never defined, and there is no field over which it is a vector
  space. Adopt composition length and prove additivity
  (Hazard~\ref{haz:step2-linear-algebra}).
\item[\textbf{C7}] Two inequivalent definitions of ``linked instance'' coexist:
  global connectivity of the value graph, in Informal Claim (IC3); and the
  per-variable condition in \S3.1. The second does not imply the first. Use the
  formal one and carry non-fragmentation separately
  (Hazard~\ref{haz:linked}).
\item[\textbf{C8}] $B/\sigma$ is the image of $B$ in $A/\sigma$, not the
  quotient of $B$; and $(C_{1}\cap\dots\cap C_{t})/\delta=\bigcap C_{i}/\delta$
  is false in general and is proved on the fly inside
  Lemma~\ref{lem:propagation}(fm) (Hazard~\ref{haz:quotient-rel}).
\item[\textbf{C9}] Three sentences beginning ``Notice that'' are used as
  lemmas and never proved: that the relation $S$ of
  Definition~\ref{def:linear-congruence}(c) is a bridge with
  $\bcol S=\proj_{1,2}(S)=\proj_{3,4}(S)=\cover\sigma$; that $\sigma$ is
  irreducible/linear/PC exactly when $0_{\mathbf A/\sigma}$ is; and that the two
  formulations of $\TS$-freeness agree. Add
  Lemma~\ref{lem:rect-basics} to the same list.
\item[\textbf{C10}] One citation silently drops a case: the source's use of
  \cite[Thm.~6.15]{ZhukStrong} omits a third alternative --- a nontrivial
  \emph{projective} subuniverse --- present in the original. Its elimination is
  legitimate under the standing Taylor hypothesis, but the source does not say
  so, and a formalization must state it as a hidden import and check it against
  what the predecessor development actually proves.

\end{enumerate}

\begin{remark}[A retracted item, and what replaced it]\label{rem:d10-retracted}
An earlier draft of this section carried an eleventh blocking item: that the
recursion-depth bound $|A|+|\Gamma|$ of \cite[Lem.~5.2]{ZhukJACM} was ``not
established for the \textsc{CheckTuple} $\to$ \textsc{SolveNonlinked} path'',
and that ``without the repair the algorithm is not polynomial''. That was
wrong, and it is withdrawn.

Lemma~5.2 bounds the depth by arguing that every recursion site except
\textsc{CheckWeakerInstance} ``reduce[s] all domains of size greater than 1''.
At the \textsc{SolveNonlinked} site this needs the instance to be neither
linked nor fragmented --- otherwise the linked-component congruence can be
trivial on some domain, that domain does not shrink, and the bound degrades
from the constant $|A|$ to $O(|A|\cdot n)$, which is fatal, since the total work
is $\mathrm{poly}(N)^{\text{depth}}$. The source does state the hypothesis and
does discharge it. It is simply not in one place:
\begin{itemize}\tightlist
\item the precondition ``an instance that is not linked and not fragmented''
  is in the prose introducing the function, and \emph{not} in the
  \texttt{Input:} line of its pseudocode;
\item it is discharged at the call site by ``$\Theta_{0}$ is not fragmented
  because it is crucial'' --- itself a one-line argument the source does not
  give: a fragmented unsatisfiable instance has an unsatisfiable part, and no
  constraint of the \emph{other} part is then crucial;
\item the step from ``not fragmented'' to ``\emph{every} domain splits'' ---
  propagate the partition along a path, which exists by non-fragmentedness and
  is value-preserving by $1$-consistency --- appears once, inside the proof of a
  \emph{different} lemma, and is never cited here.
\end{itemize}
So the mathematics is present and the theorem is not in doubt. What is absent is
any single place where the six recursion sites of Lemma~5.2 are discharged, and
that is a blueprint's job rather than a defect of the source. The same pattern
holds for the $\Gamma$-order argument, whose case~(3) needs the projections to
be unchanged under weakening, which holds because the instance is
$1$-consistent when \textsc{CheckWeakerInstance} runs --- true, and unstated.

\emph{What does remain open}, and is now a sharp question rather than a flag:
\textsc{CheckTuple} calls \textsc{SolveNonlinked} not on $\Theta_{0}$ but on
$\Theta_{0}\wedge(x_{1}\in D_{1}')\wedge\dots\wedge(x_{n}\in D_{n}')$, and
restricting domains can only \emph{destroy} non-linkedness. Each $D_{i}'$ is a
union of $\sigma_{i}$-blocks, where $\sigma_{i}$ is the minimal linear
congruence; write $\tau$ for the linked-component congruence. If
$\tau\subseteq\sigma$ then the restriction deletes whole $\tau$-blocks and
non-linkedness survives. Whether $\tau\subseteq\sigma$ at this point is not
stated in the source, and I have not checked it. That is the next thing to
verify, and it is the kind of question a blueprint exists to surface.
\end{remark}

\begin{remark}[What this list is evidence of]\label{rem:defects-meaning}
Nine blocking items in a fifty-page paper is not a high rate; it is roughly what
one should expect, and it is the reason blueprints exist. Note the
distribution: one is a genuine misstatement (D1), one is an omission (D2), two
are gaps in the mathematics that need new arguments (D3, D4), three are
hypotheses that are too weak or missing (D5, D6, D8), and two are artefacts of
the interface with the older paper (D7, D9). None of them threatens the
theorem. All of them cost time.

Remark~\ref{rem:d10-retracted} is the counterweight, and the more instructive
entry. A tenth item was listed, survived a review round, and was published ---
and was wrong, because it was taken from a reader who worked from the pseudocode
while the hypothesis it needed was in the surrounding prose. A list like this
one is exactly as reliable as the reading behind each entry, and the failure
mode is not carelessness but reading the formal-looking part of a paper and not
the sentence before it.
\end{remark}
